\chapter{Adelic Factorization Axioms}
\label{v4-arith:costello-gwilliam}

An adelic factorization construction must specify both its local
coefficients and the spaces on which positions collide. The set of
places supplies an index for the coefficients. A holomorphic curve,
a formal disc or a derived arithmetic intersection supplies different
position data. The elementary distinctions of
Section~\ref{v4-sec:word-arithmetic-positions} are the starting point.

\section{The Ran Space of \(\overline{\mathrm{Spec}\,\mathbb{Z}}\)}
\label{v4-arith:cg:ran-space}

\begin{observation}[naive Ran is terminal]
\label{v4-arith:obs:ran-naive-terminal}
The functor of nonempty finite subsets of \(X(S)\), applied to
\(X=\operatorname{Spec}\mathbb Z\) and arbitrary test schemes \(S\),
is terminal. Indeed,
\begin{equation*}
S\;\mapsto\;\mathrm{Spec}\,\mathbb{Z}(S)\;=\;\mathrm{Hom}_{\mathrm{Ring}}(\mathbb{Z},\Gamma(S,\mathcal{O}_{S}))\;=\;\{*\}
\end{equation*}
(singleton). Its only nonempty finite subset is itself. The maps
\(\operatorname{Spec}\mathbb F_p\to\operatorname{Spec}\mathbb Z\)
are points for different test schemes. They do not give several
sections over a fixed \(S\).
\end{observation}

\begin{definition}[Discrete configurations of places]
\label{v4-arith:def:ran-Arakelov}
Let \(V=\mathbb P\sqcup\{\infty\}\), where \(\mathbb P\) is the set
of rational primes. The discrete place configurations are
\(\operatorname{Fin}_{+}(V)\). When multiplicities are retained,
use \(\operatorname{Mult}_{+}(V)=\coprod_{n\ge1}V^n/\mathfrak S_n\).
These are combinatorial index sets. They specify neither a scheme
structure on an Arakelov compactification nor a normal neighbourhood
of a position diagonal. We write \(\overline C\) for the arithmetic
curve only when its additional geometric data have been specified.
\end{definition}

\begin{remark}
An order on \(V\), or the multiplicity of a place in a multiset,
does not produce a normal coordinate. Such coordinates belong to a
separately chosen local position space. The cut construction of
Chapter~\ref{v4-ch:words-collisions-units} uses a smooth curve and
its differential modules; its hypotheses are additional to the
choice of \(V\).
\end{remark}

\section{Factorization for Disjoint Place Labels}
\label{v4-arith:cg:factorization-iso-trivial}

\begin{observation}
\label{v4-arith:obs:factorization-iso-trivial}
Choose a common commutative coefficient ring \(A\) and complexes
\(H_v\in\mathcal D(A)\). Set
\(F(S)=\bigotimes_{v\in S,A}^{\mathbf L}H_v\) for a finite set
of places. For disjoint sets \(S,T\), the tensor associator gives
\(F(S\sqcup T)\simeq F(S)\otimes_A^{\mathbf L}F(T)\).
Its coherences follow from those of derived tensor, as in
\eqref{v4-eq:word-discrete-factorization}.
\end{observation}

\begin{corollary}[Scope of discrete factorization]
\label{v4-arith:cor:no-nontrivial-factorization}
The preceding construction gives a coherent factorization for
disjoint finite subsets of \(V\). It requires no maps at coincident
positions. Identifying it with a geometric arithmetic factorization
system requires the position spaces, their collision coefficients
and comparison maps in addition to the disjoint tensor identities.
\end{corollary}
\begin{proof}
Every map used in the construction is an associator or symmetry of
the finite coefficient tensor. Its source and target are indexed
by disjoint sets. No diagonal coefficient occurs among these maps.
\end{proof}

\section{The Diagonal and \(\mathbb{Z}_{p}\,\widehat{\otimes}\,\mathbb{Z}_{p}\)}
\label{v4-arith:cg:diagonal-Zp-tensor}

\begin{observation}
\label{v4-arith:obs:Zp-tensor-collapses}
For the ordinary \(p\)-adic completion of the indicated tensor,
\begin{equation}
\mathbb{Z}_{p}\,\widehat{\otimes}_{\mathbb{Z}}\,\mathbb{Z}_{p}\;\cong\;\mathbb{Z}_{p},
\end{equation}
where the left side means
\(\varprojlim_r(\mathbb Z_p\otimes_{\mathbb Z}\mathbb Z_p)/p^r\).
Multiplication is an isomorphism modulo \(p^r\) for every \(r\),
since both factors then give \(\mathbb Z/p^r\). The inverse limit
proves the assertion. This computation concerns the coefficient
base, with no position variable.
\end{observation}

\begin{remark}
A chosen formal position disc has self-product
\(\mathbb Z_p[[t_1,t_2]]\) and nonzero diagonal ideal
\((t_1-t_2)\). The derived closed-point intersection instead has
the normal class in
\(\mathbb F_p\otimes_{\mathbb Z}^{\mathbf L}\mathbb F_p
 \simeq\mathbb F_p\oplus\mathbb F_p[1]\).
These are distinct coefficient models, as proved in
Proposition~\ref{v4-prop:word-finite-place-models}.
Crystalline or prismatic data can be added when a construction
requires their Frobenius structure; the completed base tensor
does not force that choice.
\end{remark}

\begin{theorem}[Coefficient completion and position diagonals]
\label{v4-arith:thm:arch-vs-finite-asymmetry}
The diagonal of a complex position line has normal coordinate
\(z_1-z_2\). A formal position disc over \(\mathbb Z_p\) has
normal coordinate \(t_1-t_2\). Both ideals are nonzero, although
the ordinary completed base tensor of
Observation~\ref{v4-arith:obs:Zp-tensor-collapses} is \(\mathbb Z_p\).
\end{theorem}
\begin{proof}
In both position rings, change variables to the second coordinate
and the difference of the two coordinates. Restriction to the
diagonal evaluates the difference at zero and has the displayed
nonzero kernel. The base-tensor assertion is the observation above.
\end{proof}

\section{Pointed Data at Each Finite Place}
\label{v4-arith:cg:E0-finite-places}

\begin{remark}[pointed local data and factorization]
\label{v4-arith:prop:locally-constant-discrete}
An \(E_0\)-algebra in vector spaces is a pointed vector space.
The vacuum therefore equips each polynomial oscillator module with such
an underlying datum. This does not construct a factorization algebra on
\(\mathrm{Spf}\,\mathbb Z_p\). Classification results for locally constant
factorization on manifolds require their geometric hypotheses; a formal
scheme is not made a zero-dimensional manifold by this notation.
\end{remark}

\begin{definition}[minimal adelic factorization structure]
\label{v4-arith:def:minimal-adelic-factorization}
The proposed adelic factorization datum consists of:
\begin{itemize}
\item At \(\infty\): genuine holomorphic factorization algebra on
\(\mathbb{C}_{\infty}\) in the CG sense (Vol.~I~Def.~3.1.4), with
Heisenberg OPE \(j(z)j(w)\sim k_{\infty}/(z-w)^{2}\), \(k_{\infty}=1\).
\item At each finite \(p\): locally constant prestack on
\(\mathrm{Spf}\,\mathbb{Z}_{p}\) with constant stalk \(\mathcal{H}^{(p)}\)
and a specified Adams endomorphism \(\psi_p\). Its underlying
pointed module is an \(E_0\)-datum; its geometric structure maps are
additional requirements.
\item Gluing data: restricted tensor product over spherical vectors
\(|0\rangle_{p}\) and explicit local-to-global comparison maps. An automorphic
representation is further data, not a consequence of the tensor notation.
\end{itemize}
\end{definition}

\section{The Holomorphic-Topological Structure}
\label{v4-arith:cg:HT-structure}

\begin{remark}[Costello HT factorization, arXiv:1303.2632; Costello 2020]
Costello's holomorphic-topological factorization algebras live on
\(\mathbb{C}^{n_{H}}\times\mathbb{R}^{n_{T}}\) with \(n_{H}\) holomorphic
and \(n_{T}\) topological directions. The axioms blend holomorphic
OPE (on the \(\mathbb{C}^{n_{H}}\) factor) and locally constant
(\(E_{n_{T}}\)-algebra structure on the \(\mathbb{R}^{n_{T}}\) factor).
\end{remark}

\begin{conjecture}[adelic HT structure]
\label{v4-arith:thm:adelic-HT}
\emph{Conjectural statement.} On \(\overline{C}\) the natural HT
structure is
\begin{equation*}
\underbrace{\mathbb{C}_{\infty}}_{\text{holomorphic}}
\;\times\;\underbrace{\mathrm{Spf}\,\widehat{\mathbb{Z}}}_{\text{adelic / }E_{0}\text{ at each place}},
\end{equation*}
with the profinite completion \(\widehat{\mathbb{Z}}=\prod_{p}\mathbb{Z}_{p}\)
totally disconnected. Locally constant factorization on
\(\mathrm{Spf}\,\widehat{\mathbb{Z}}\) decomposes as a product of
\(E_{0}\)-stalks. The non-trivial chiral content concentrates at
\(\infty\).
\end{conjecture}

\begin{definition}[relative factorization window with adelic base]
\label{v4-arith:def:relative-adelic-factorization}
A relative factorization window for
\(\mathcal{V}^{\mathrm{prim}}\), when supplied, is a holomorphic
factorization algebra on
\(\mathbb{C}_{\infty}\) \emph{over the base}
\(\mathrm{Spec}\,\widehat{\mathbb{Z}}\), with the base acting by
Adams operations \(\{\psi_{p}\}_{p}\) and Hecke operators
\(\{T_{p}\}_{p}\):
\begin{equation}
\mathcal{F}_{\mathrm{CG}}(\mathcal{V}^{\mathrm{prim}})
\;\in\;\mathrm{FactAlg}_{\mathbb{C}_{\infty}}\bigl(\mathrm{Mod}_{\widehat{\mathbb{Z}},\psi}\bigr),
\end{equation}
where the coefficient category and its scalar comparison at infinity
must be specified. The symbol
\(\mathrm{Mod}_{\widehat{\mathbb Z},\psi}\) records modules with a
named compatible commuting action; it does not identify divisor
substitutions with Adams maps. The full divisor family fails that
commutation requirement by Equation~\ref{v4-arith:ado:boundary-commutator}.
\end{definition}

\section{Five Axioms (F1)--(F5) for the Minimal Adelic Factorization}
\label{v4-arith:cg:five-axioms}

\begin{definition}[Costello--Gwilliam arithmetic factorization]
\label{v4-arith:def:CG-arithmetic-factorization}
The following are requirements on an arithmetic factorization datum,
not consequences of the polynomial oscillator representation:

\paragraph{(F1) Archimedean factorization.} Holomorphic
factorization algebra on \(\mathbb{C}_{\infty}\) in the sense of
CG~Vol.~I~Def.~3.1.4, with OPE
\(j(z)j(w)\sim k_{\infty}/(z-w)^{2}\), \(k_{\infty}=1\).

\paragraph{(F2) Coefficient category.} Values in
\(\mathrm{Mod}_{\widehat{\mathbb{Z}},\psi}\):
\(\widehat{\mathbb{Z}}\)-modules with commuting Adams operations
\(\psi_{p}\colon\mathcal{M}\to\mathcal{M}\) satisfying
\(\psi_{p}\psi_{q}=\psi_{q}\psi_{p}\) and \(\psi_{p}|_{\mathbb{Z}_{p}}=\mathrm{id}\).

\paragraph{(F3) Hecke-eigen decomposition.} Global sections
decompose as a restricted tensor product
\begin{equation*}
\mathcal{V}^{\mathrm{prim}}(\mathbb{C}_{\infty})\;=\;\bigotimes_{p}^{\prime}\mathcal{H}^{(p)}
\end{equation*}
with specified reference vectors \(|0\rangle_p\) and separately
specified intertwiners for its Adams and Hecke operators.

\paragraph{(F4) Archimedean--finite compatibility.} The
archimedean Heisenberg level is \(k_{\infty}=1\), hence
\(v_{p}(k_{\infty})=0\) for every finite prime. The finite
arithmetic datum is carried by the Satake eigenvalue \(a_{p}\), not by
a \(p\)-adic valuation of the archimedean level.

\paragraph{(F5) Regular \(p\)-adic product.} At each finite place
\(p\), the local factor \(\mathcal{H}^{(p)}\) carries regular
(non-singular) commutative multiplication and Adams endomorphism
\(\psi_p\) among its specified operations. The product
\(\mathcal{H}^{(p)}\otimes\mathcal{H}^{(p)}\to\mathcal{H}^{(p)}\) is
the commutative multiplication in the Fock module: a commutative
algebra structure, distinct from a BD factorization-algebra fusion.
\end{definition}

\begin{remark}[scope of the five conditions]
\label{v4-arith:thm:CG-conclusion}
Conditions (F1)--(F5) specify proposed local objects, coefficient actions
and comparisons. The commutative product and index-shift action on the
finite polynomial carrier satisfy the algebraic part of (F5).
The geometric structure maps in (F1), the coefficient comparison in (F2),
and the global tensor and operator identifications in (F3) require
separate constructions. No arithmetic factorization existence theorem
follows from this local polynomial calculation.
\end{remark}

\section{Classical Factorization at Finite Places: Absence of a BV Action Functional}
\label{v4-arith:cg:BV-quantization-fails}

\begin{remark}[BV quantization requires action functional]
CG~Vol.~I~\S4 constructs a factorization algebra of observables as
a BV-quantization of a classical field theory \((\mathcal{E},S)\).
For the free boson on \(C/\mathbb{C}\):
\(\mathcal{E}=\Gamma(C,\Omega^{1}_{C})[-1]\), action
\(S[\phi]=\int_{C}\phi\wedge\bar{\partial}\phi\), the
quantum factorization algebra is the Heisenberg factorization algebra
(CG~Vol.~I~Thm~4.6.1; as a translation-equivariant factorization
algebra on \(\mathbb{C}\), this recovers the Heisenberg vertex algebra,
but the two structures are distinct: factorization algebra
\(\neq\) vertex algebra).
\end{remark}

\begin{observation}[arithmetic differential forms and the free-field obstruction]
\label{v4-arith:obs:no-classical-arith-FT}
The sheaf of Kähler differentials \(\Omega^{1}_{\mathbb{Z}/\mathbb{Z}}\)
on \(\overline{\mathrm{Spec}\,\mathbb{Z}}\) is zero in the ordinary
algebraic sense; Arakelov geometry replaces it with a metrised line
bundle at the archimedean place. A classical free-field action at a
finite place \(p\) would require a section of
\(\Omega^{1}_{\overline{C}/\mathbb{F}_{1}}\) in a neighbourhood of
\(p\); the requisite \(\mathbb{F}_{1}\)-geometry remains
unconstructed.
\end{observation}

\begin{remark}
Borger's Frobenius-lift criterion has the torsion-free scope stated in
Theorem~\ref{v4-arith:thm:wilkerson}. The monomial and Hall integral lattices
carry different reductions. The shifts \(\psi_n\) preserve both, but
are Frobenius lifts only on the latter. The divisor maps \(T_n\) also
preserve both lattices and fail that condition even on the Hall lattice.
These operations supply no action functional, BV bracket or geometric
factorization map.
\end{remark}

\begin{theorem}[finite polynomial Adams datum]
\label{v4-arith:thm:classical-finite-places}
For each finite arithmetic place \(p\), let
\(P_{\mathbb Z_p}=\mathbb Z_p[x_k:k\ge1]\). Its ordinary product,
vacuum \(1\), and substitutions \(\psi_n(x_k)=x_{nk}\) give a unital
commutative algebra with an \(\mathbb N^{\times}\)-action.
The prescribed \(\psi_p\) is not a Frobenius lift on this monomial
lattice. An action functional or factorization algebra on
\(\mathrm{Spf}\,\mathbb Z_p\) is additional data.
\end{theorem}
\begin{proof}
The polynomial product and substitutions are defined over
\(\mathbb Z_p\), and \(\psi_m\psi_n(x_k)=x_{mnk}\) proves the monoid
law on generators. The two distinct monomials in
\(\psi_p(x_1)-x_1^p=x_p-x_1^p\) remain distinct modulo \(p\).
This proves the claimed failure of this particular Frobenius lift;
it proves no impossibility of other integral structures or field theories.
\end{proof}

\section{Factorization Boundary}
\label{v4-arith:cg:summary}

\begin{enumerate}
\item The finite-subset functor of
\(\operatorname{Spec}\mathbb Z(S)\) is terminal. Finite sets and
multisets of places are separate discrete index sets
(Observation~\ref{v4-arith:obs:ran-naive-terminal},
Definition~\ref{v4-arith:def:ran-Arakelov}).
\item Finite coefficient tensors give coherent factorization
for disjoint place labels. Their geometric collision maps require
additional data
(Corollary~\ref{v4-arith:cor:no-nontrivial-factorization}).
\item The completed tensor product
\(\mathbb{Z}_{p}\,\widehat{\otimes}_{\mathbb{Z}}\,\mathbb{Z}_{p}\cong\mathbb{Z}_{p}\)
makes the completed base diagonal an isomorphism. Formal position
discs and derived closed-point intersections retain normal data
(Observation~\ref{v4-arith:obs:Zp-tensor-collapses},
Theorem~\ref{v4-arith:thm:arch-vs-finite-asymmetry}).
\item At each finite place the vacuum gives a pointed module.
A factorization structure on a formal scheme is additional geometry
(Remark~\ref{v4-arith:prop:locally-constant-discrete}).
\item The proposed relative datum includes a
holomorphic factorization algebra on \(\mathbb{C}_{\infty}\) with
\(\widehat{\mathbb{Z}}\)-module coefficients and commuting Adams
operations, as in (F1)--(F5)
(Definition~\ref{v4-arith:def:relative-adelic-factorization},
Definition~\ref{v4-arith:def:CG-arithmetic-factorization}).
\item The monomial finite-place carrier has an Adams action whose
prime shift fails the Frobenius congruence. BV quantization requires a
separately supplied classical action functional
(Theorem~\ref{v4-arith:thm:classical-finite-places}).
\end{enumerate}

A Costello--Gwilliam window, if constructed with the stated data, is a
holomorphic factorization algebra on \(\mathbb{C}_{\infty}\) with an
adelic local system decoration (axioms (F1)--(F5)); the full Hochschild
trace class itself is not thereby made a vertex algebra.
